\section{Vehicle model}
The vehicle model must capture the main relationship between steering, longitudinal motion, and the resulting trajectory while remaining sufficiently simple for the numerical optimization problem. Since this project focuses on the minimum-lap-time optimization problem rather than on detailed vehicle dynamics, we use a kinematic bicycle model.

The kinematic bicycle model describes the vehicle motion using only a few states and control inputs. It is simple enough for numerical optimization while still capturing the main effects of steering and acceleration. More detailed effects, such as tire slip, load transfer, suspension motion, and transient tire forces, are neglected.
\subsection{The bicycle model}\;\\
The commonly used kinematic bicycle model is a simplification of the dynamics of a car. It represents the two wheels of each axle by a single equivalent wheel, resulting in one front wheel and one rear wheel. Only the front wheel is used for steering, and tire slip is neglected.

The model is described by the position of the car in global coordinates $(x,y)$, the velocity $v$, the heading $\psi$, and the steering angle $\delta$. The dynamics of the model are given by
\begin{align*}
    \dot{x} &= v \cos(\psi) \\
    \dot{y} &= v \sin(\psi) \\
    \dot{\psi} &= \frac{v}{L} \tan(\delta) \\
    \dot{v} &= a
\end{align*}
The parameter $L$ denotes the wheelbase of the car, i.e. the distance between the front and rear wheel, and the acceleration of the car respectively. For more details on this model, see \cite{Rajamani2012}.

However, for driving around a given track it is inconvenient to use global coordinates because it is more complicate to formulate the constraints. Additionally, we only need the position on the track and do not care for the absolute position in global coordinates.

Therefore, we want to formulate the model using track coordinates. The centerline distance traveled will be given as $s$ and $e_y$ denotes the distance from the centerline, $e_\psi$ denotes the angle between the cars heading and the tangent of the centerline and $v$ still denotes the velocity of the vehicle.

The position of the vehicle can be expressed relative to the track centerline as
\begin{align*}
    \mathbf{p}(t)
    =
    \mathbf{p}_c(s(t))
    +
    e_y(t)\mathbf{N}(s(t)),
\end{align*}
where $\mathbf{p}_c(s)$ denotes the centerline position and
$\mathbf{N}(s)$ is its unit normal vector.

Differentiating with respect to time gives
\begin{align*}
    \dot{\mathbf{p}}
    &=
    \frac{d\mathbf{p}_c}{ds}\dot{s}
    +
    \dot{e}_y\mathbf{N}
    +
    e_y\frac{d\mathbf{N}}{ds}\dot{s}.
\end{align*}
Using
\begin{align*}
    \frac{d\mathbf{p}_c}{ds}=\mathbf{T},
    \qquad
    \frac{d\mathbf{N}}{ds}=-\kappa(s)\mathbf{T},
\end{align*}
where $\mathbf{T}(s)$ is the unit tangent vector, we obtain
\begin{align*}
    \dot{\mathbf{p}}
    =
    \left(1-\kappa(s)e_y\right)\dot{s}\,\mathbf{T}
    +
    \dot{e}_y\,\mathbf{N}.
\end{align*}

The vehicle velocity can also be decomposed into components tangent and
normal to the centerline:
\begin{align*}
    \dot{\mathbf{p}}
    =
    v\cos(e_\psi)\mathbf{T}
    +
    v\sin(e_\psi)\mathbf{N}.
\end{align*}

Equating the tangential and normal components gives
\begin{align*}
    \left(1-\kappa(s)e_y\right)\dot{s}
    &=
    v\cos(e_\psi), \\
    \dot{e}_y
    &=
    v\sin(e_\psi).
\end{align*}
Therefore,
\begin{align*}
    \dot{s}
    =
    \frac{v\cos(e_\psi)}
    {1-\kappa(s)e_y}.
\end{align*}
This relation is commonly used in Frenet-coordinate formulations of vehicle
dynamics \cite{Reiter2023Frenet}.

The heading error is defined as
\begin{align*}
    e_\psi = \psi-\psi_c(s),
\end{align*}
where $\psi_c(s)$ denotes the heading of the track centerline.

Differentiating with respect to time and applying the chain rule yields
\begin{align*}
    \dot e_\psi
    =
    \dot\psi
    -
    \frac{d\psi_c}{ds}\dot s.
\end{align*}

The cars own heading with respect to the tangent of the centerline changes because of steering and the heading of the track changes because of the curvature of the track.

Since the derivative of the centerline heading with respect to the arc length is equal to the local track curvature,
\begin{align*}
    \frac{d\psi_c}{ds}=\kappa(s),
\end{align*}
the heading-error dynamics become

\begin{align*}
    \dot e_\psi = \frac{v}{L} \tan(\delta) - \kappa(s)\dot s.
\end{align*}

This results in the following vehicle model expressed in Frenet coordinates.
The state vector is defined as
\begin{align*}
    z(t)
    =
    \begin{bmatrix}
        s(t)\\
        e_y(t)\\
        e_\psi(t)\\
        v(t)
    \end{bmatrix},
\end{align*}
where \(s(t)\) denotes the progress along the track centerline,
\(e_y(t)\) the lateral offset from the centerline,
\(e_\psi(t)\) the heading error, and \(v(t)\) the vehicle velocity.

The control vector is given by
\begin{align*}
    u(t)
    =
    \begin{bmatrix}
        a(t)\\
        \delta(t)
    \end{bmatrix},
\end{align*}
where \(a(t)\) is the longitudinal acceleration and \(\delta(t)\) is the
steering angle.

The state dynamics can then be written compactly as
\begin{align*}
    \dot z(t)
    =
    f\bigl(z(t),u(t)\bigr)
    =
    \begin{bmatrix}
        \dfrac{v(t)\cos(e_\psi(t))}
        {1-\kappa(s(t))e_y(t)}
        \\[3ex]
        v(t)\sin(e_\psi(t))
        \\[1.5ex]
        \dfrac{v(t)}{L}\tan(\delta(t))
        -
        \kappa(s(t))
        \dfrac{v(t)\cos(e_\psi(t))}
        {1-\kappa(s(t))e_y(t)}
        \\[1.5ex]
        a(t)
    \end{bmatrix}.
\end{align*}

The kinematic bicycle model neglects tire slip, load transfer,
suspension dynamics and transient tire forces.
Consequently, it is not intended to accurately describe the vehicle behaviour
at the limits of adhesion.

Nevertheless, it captures the main relationship between steering,
longitudinal acceleration and the resulting vehicle trajectory while keeping
the optimization problem sufficiently simple.